\documentclass[12pt]{article}
\usepackage[utf8]{inputenc}
\usepackage[T1]{fontenc}
\usepackage{hyperref}
\title{Yamongeri Language Resources for K-12 Curriculum Development}
\author{Compiled by Claude Sonnet 5.5 (claude-sonnet-5-5) for Andr\'as Kornai}
\date{}
\begin{document}
\maketitle

\section{Introduction}
Yamongeri (ISO 639-3: ymg; also Yamongiri) is a Bantu language of the Democratic Republic of the Congo. It is spoken in \'Equateur province, in the northwest of the country south of the Congo River. Joshua Project estimates about 29,000 people. Wikitongues lists it as Niger-Congo with an EGIDS status of 6a (vigorous). The Wikipedia article on Mongo places Yamongeri in the wider Mongo cluster. The search results gave no evidence of any school materials, a standard orthography or a published curriculum in Yamongeri. This report is therefore short.

\section{Primary Educational Resources}
No primers, textbooks or school programmes in Yamongeri turned up in the searches. Joshua Project notes that formal schooling is often limited in rural \'Equateur.

\section{Linguistic and Reference Materials}
\begin{itemize}
\item Wikitongues language page for ymg: \url{https://wikitongues.org/languages/ymg/}
\item Wikipedia, Mongo language (the wider cluster, Lomongo/Nkundo): \url{https://Www.wikipedia.org/wiki/Nkundo}
\end{itemize}
These are general reference pages. No grammar or dictionary of Yamongeri itself was found.

\section{Religious and Community Resources}
Joshua Project reports that Bible translation has started, with no audio Bible available. It describes the people group as partially reached, with Christianity as the main religion.
\begin{itemize}
\item Joshua Project, Yamongeri people group: \url{https://joshuaproject.net/people_groups/15977}
\item Joshua Project, language page: \url{https://joshuaproject.net/languages/ymg}
\end{itemize}

\section{Digital and Audio Resources}
No audio, video or software resources specific to Yamongeri were found. Joshua Project lists no audio Bible.

\section{Curriculum Development Context}
Very little is documented. The written status of the language is listed as unknown, and no school use is known. Any K-12 curriculum would need new fieldwork, a decision on orthography and community involvement. Related Mongo (Lomongo) materials might serve as a starting point, but they should not be assumed to match Yamongeri.

\end{document}